\documentclass[11pt]{article}
% preamble.tex — shared preamble for the Flyxion connective-essay series
% Compile target: XeLaTeX. \input this file, or copy its contents, at the
% top of each essay after \documentclass.

\usepackage{fontspec}
\setmainfont{TeX Gyre Termes} % Times-like serif; falls back gracefully if absent
\usepackage{amsmath,amssymb,amsthm}
\usepackage[margin=1.25in]{geometry}
\usepackage[hidelinks]{hyperref}
\usepackage{microtype}
\usepackage{enumitem}
\usepackage{footmisc}

\newtheorem{claim}{Claim}
\newtheorem{definition}{Definition}
\newtheorem{observation}{Observation}

% --- Corpus vocabulary macros -------------------------------------------
\newcommand{\admissible}{\textsc{admissible}}
\newcommand{\inadmissible}{\textsc{inadmissible}}
\newcommand{\repair}{\textit{repair}}
\newcommand{\continuation}{\textit{continuation}}
\newcommand{\rsvp}{\textsc{rsvp}}

% --- Title block helper ---------------------------------------------------
% Usage: \essaytitle{Title}{Subtitle or tagline}
\newcommand{\essaytitle}[2]{%
  \title{#1\\[0.3em]\large #2}
  \author{Flyxion\\\normalsize Independent Researcher}
  \date{\today}
}

\pagestyle{plain}

\essaytitle{Aviation Incident Ledgers and the Repair of Safety Knowledge}{ASRS, NTSB, and a working counterexample to ledger skepticism}
\begin{document}
\maketitle
\begin{abstract}
Debates about incident ledgers in emerging domains often assume that a ledger without a clean denominator, compulsory reporting, or independent observations cannot sustain serious inference. Aviation safety offers a mature counterexample. NASA's Aviation Safety Reporting System (ASRS) and the National Transportation Safety Board (NTSB) do not solve every measurement problem, and the field itself is candid about underreporting and selection effects. But they do show that a reporting ledger can become action-guiding without first becoming an ideal statistical sample. Aviation does this by treating near-misses as structured evidence, by coding reports for severity and contributing factors, by combining multiple reports on a single event, and by connecting findings to recommendations, airworthiness directives, training changes, and procedure revisions. The result is not a naive count-book of isolated incidents. It is a continuing evidentiary apparatus in which repeated, non-independent disturbances accumulate weight precisely because they disclose common-mode failure. That mature practice undercuts a recurring premise in critiques of AI incident ledgers: the premise that imperfect denominator knowledge makes the ledger itself epistemically unrepairable.
\end{abstract}

\section{Introduction}
The modern aviation safety apparatus was built around a lesson that other high-risk fields keep rediscovering: waiting for disasters alone is too expensive. Since 1976, the ASRS has collected voluntary, confidential reports from pilots, controllers, mechanics, flight attendants, dispatchers, and others.\footnote{NASA, ``ASRS Program Briefing'' and ``Summary,'' at the ASRS overview pages; FAA Advisory Circular 00-46F, \textit{Aviation Safety Reporting Program}, sections 6--10. The FAA instituted the reporting program in 1975, and NASA began operating ASRS under the FAA--NASA agreement in 1976.} In recent years that intake has been on the order of roughly 100,000 reports per year, so the epistemic problem is not how to dignify a lone anecdote but how to code and triage a very large voluntary stream. NASA administers it as a neutral third party rather than the FAA, specifically to encourage disclosure that would otherwise be chilled by fear of punishment. Reports are de-identified before entry into the public database, and the system publishes recurring lessons through outlets such as \textit{CALLBACK}.\footnote{NASA ASRS, ``Confidentiality and Incentives to Report''; ``About ASRS Data''; and ``ASRS Coding Taxonomy.'' NASA states that names are removed, dates and times may be generalized or eliminated, and that no reporter's identity has been breached under NASA management.} In parallel, the NTSB investigates accidents as an independent safety body whose work is organized around factual reconstruction, probable-cause analysis, and safety recommendations rather than ordinary enforcement or tort liability.\footnote{NTSB, ``The Investigative Process''; 49 U.S.C. \S 1154(b) (restricting the use of NTSB reports in civil actions for damages).}

This matters because some critiques of incident ledgers in newer fields, especially AI safety, focus on three familiar worries. First, if one does not know the denominator --- how many opportunities for failure existed --- raw incident counts are said to be nearly meaningless. Second, near-miss reports are said to be too voluntary, selective, or anecdotal to support disciplined learning. Third, incidents are said to be too non-independent for aggregation: one laboratory, one model family, or one deployment culture can generate clusters that look like multiple events while really expressing only a single hidden cause. Those worries are real. They are also, taken as objections to the very possibility of a useful ledger, too strong. Aviation has spent decades learning from exactly such data.

The point is not that aviation has perfect numbers. It does not. The point is that it has a working method for turning imperfect, biased, and interdependent reports into \continuation-worthy safety knowledge.

\subsection{Step 2: The unexamined premise}
The unexamined premise behind a certain class of ledger skepticism is that an incident ledger is fundamentally a defective statistics table. On that view, unless one can attach each entry to a clean exposure measure and defend something like independent sampling, the ledger cannot license more than story-telling. A voluntary near-miss report then appears as a bias to be regretted, not as a form of signal to be structured.

Aviation shows a different possibility. A safety ledger need not begin by answering the global question ``how many events per unit exposure occurred in the whole domain?'' It can begin by asking a narrower and often more tractable question: which reports are \admissible as safety signal, how should they be coded, what latent conditions do they implicate, and what corrective action would make recurrence less likely? That is already enough to support learning, prioritization, and intervention.

\subsection{Step 3: Temporal versus constitutive priority}
There are two different claims one might make about the relation between counts and safety judgment. The first is temporal: sometimes an analyst obviously wants rates later, once enough data and exposure estimates exist. The second is constitutive: one might claim that no meaningful safety inference exists \emph{until} a clean denominator is in hand. The present essay denies the second claim, not the first.

In aviation, denominator information often arrives late, partially, or at the wrong level of granularity. Flight hours, sectors, fleet types, weather conditions, training environments, maintenance regimes, and airport layouts all matter, and the relevant exposure unit varies by question. Yet the field does not suspend judgment until every denominator is globally normalized. Constitutively, the load-bearing work is done earlier by filters of seriousness, coding, corroboration, and repairability. A report can be good enough to change training, phraseology, maintenance checks, or cockpit procedure well before it supports a satisfying population-rate estimate.

\section{Minimum apparatus: admissibility and repair}
Only two pieces of the larger vocabulary are needed here.

\begin{definition}
An aviation safety report is \admissible, for present purposes, when it can be entered into a shared learning process without pretending to a precision it does not possess: it is identifiable as a safety-relevant event or condition, de-identified enough to be reportable, classifiable enough to be compared with others, and concrete enough to support at least a candidate intervention.\footnote{An adjacent empirical result makes the same point quantitatively. Xinyu Pang et al., \emph{DNative-Twin: Decision Graphs and Digital Twins for Reconstructable Agentic Decisions}, arXiv:2609.03787 (2026), \url{https://arxiv.org/abs/2609.03787}, find that recall for an injected, unresolved failure rises from 0 with bare causal-graph structure alone, to 0.667 once replay-contract state is added, to 1.0 once verification evidence is additionally available. Retained history is not automatically diagnostically sufficient history; the same distinction is at stake in the difference between a raw incident count and a coded, corroborated ASRS report.}
\end{definition}

A \repair{} is then not merely the correction of one damaged aircraft or one erroneous clearance. It is any loop-closing intervention --- revised procedure, training change, alerting message, airworthiness directive, checklist redesign, equipment modification, or recommendation tracking --- that lets future operations continue without reintroducing the same contradiction under a new disguise. That use follows the narrow claim made in \textit{Distinction Holonomy}: repair need not restore a prior state byte for byte; it must preserve a path-consistent way of going on.

A concrete aviation example is the rise of Crew Resource Management after late-1970s accidents such as United Airlines Flight 173, where a crew fixated on a landing-gear problem exhausted its fuel while cockpit communication and challenge norms failed, reinforced by the broader post-Tenerife push to treat cockpit coordination, authority gradients, and cross-checking as safety-critical. The repair was not a restored airframe. It was a change in admissible cockpit practice: explicit training in communication, workload sharing, and junior-crew challenge behavior so that later flights could continue without reproducing the same organizational failure mode.

\section{ASRS: voluntary near-miss reporting as structured signal}
The clearest counterexample to ledger skepticism is the ordinary processing of an ASRS submission. NASA's own descriptions emphasize several linked design choices. Reporting is voluntary. Confidentiality is strict. Names and organizational identifiers are removed. Dates, times, and related details that might permit re-identification are generalized or eliminated. For non-accident, non-criminal reports, the identifying strip is returned to the reporter and not retained in the working file. The incentive structure is equally important: the FAA has long committed not to use ASRS reports in enforcement actions within the program's limits, and in qualifying cases it waives penalties for unintentional violations. The system was built this way precisely because punitive environments suppress disclosure.

If one approached the system with a naive sampling ideal, nearly every one of these features could be redescribed as a defect. Voluntary reporting creates self-selection. De-identification sacrifices contextual detail. Immunity incentives may skew what gets disclosed. NASA does not fully investigate every report in the manner of a formal accident inquiry. ASRS itself warns that the presence of reports in the database cannot be used to infer the prevalence of a problem across the whole National Airspace System.

But that warning is not an admission of uselessness. It is an admissibility rule. The report is not being admitted as a rate estimate. It is being admitted as evidence of vulnerability. That is why ASRS codes reports into structured fields and taxonomies rather than leaving them as an archive of anecdotes. The online coding documentation distinguishes fixed fields from textual ones, and the public descriptions highlight categories such as anomaly type, contributing factors, human factors, phase of flight, weather, airspace, and outcome. Multiple reports concerning a single event can be combined into one record, with supplemental accounts used to enrich the primary one. In other words, the system does not pretend each narrative is an independent draw. It explicitly merges perspectives when they belong to one occurrence.

The \textit{CALLBACK} bulletin shows the epistemic intent of the apparatus. It republishes de-identified reports as lessons learned, often with analyst commentary or thematic grouping. The purpose is neither prosecution nor actuarial closure. The purpose is to circulate practical warning signs: unstable approaches, readback-hearback failures, runway confusion, maintenance workarounds, automation surprises, fatigue patterns, and communication ambiguities. A single report rarely proves much. A coded and recurrent family of reports can justify attention long before the field has an elegant denominator.

This is exactly where a ledger becomes a \continuation{} device rather than a passive archive. By soliciting near-misses, the system captures disturbances that ended short of catastrophe but reveal the contour of the hazard. The gain is not that near-misses are representative in the survey-statistics sense. It is that they are often mechanistically informative. A near-midair conflict, a mis-set altitude, or a maintenance confusion may be rare in the absolute and still extraordinarily rich about where the defensive layers are thinning.

\section{NTSB investigation and the evidentiary role of non-independence}
The NTSB provides the complementary half of the picture. Its general investigative process is publicly described as notification and launch, on-scene fact gathering, analysis, probable-cause determination, final report adoption, and then advocacy for safety recommendations beyond the report itself. That final step is crucial. The investigation is not complete at narrative closure. It is complete only when the findings are connected to prospective change.

Here James Reason's ``Swiss cheese'' picture remains useful, even if only as a first approximation.\footnote{James Reason, \emph{Human Error} (Cambridge: Cambridge University Press, 1990).} Its enduring contribution is not a cartoon of holes lining up, but the insistence that accidents are layered organizational events rather than isolated front-line mistakes. Investigators look not only for the active failure at the sharp end but for latent conditions --- training weakness, supervision gaps, maintenance culture, design assumptions, scheduling pressure, poor interfaces, ambiguous procedures --- that make similar breakdowns likely elsewhere.

The 2009 Colgan Air Flight 3407 crash is illustrative.\footnote{NTSB, \textit{Loss of Control on Approach, Colgan Air, Inc., Operating as Continental Connection Flight 3407, Bombardier DHC-8-400, N200WQ, Clarence Center, New York, February 12, 2009} (Aircraft Accident Report, 2010).} The NTSB's findings did not stop with the crew's inappropriate response to the stick shaker. The investigation also emphasized training deficiencies, professional standards, fatigue issues, and shortcomings in carrier and regulatory oversight. Recommendations extended outward accordingly: pilot records, stall-recovery training, fatigue management, and broader operational oversight all became part of the remedial package. This is not an analyst violating the independence assumption by illicitly generalizing from one event. It is an investigator recognizing that the event is evidentially powerful \emph{because} it exposes a shared system condition.

In that sense, non-independence is often the very thing one wants to detect. If separate incidents are linked by the same organizational weakness, then their failure to be independent does not nullify learning; it explains why accumulation matters.\footnote{A recent formal treatment of exactly this point comes from a different domain: Yifan Wu et al., \emph{Safety Does Not Compose: Non-Decaying Loop State for Autonomous LLM Agents}, arXiv:2608.27141 (2026), \url{https://arxiv.org/abs/2608.27141}, prove that when evidence of a failure is distributed across repeated cycles, any monitor that resets between cycles is information-theoretically unable to detect it, while a monitor that carries state forward across cycles can. The aviation case differs in mechanism but shares the underlying structure: what looks like several independent reports may in fact be one hazard visible only in aggregate.} Mature safety analysis therefore cross-references nominally separate events against recurring causal patterns. The ledger thickens around the common mode. When later reports resemble earlier ones in phase of flight, crew interface, maintenance handoff, or organizational setting, they do not merely increase a count. They sharpen the contour of a latent hazard.

That is why recommendations function as the field's \repair{} operator. Findings alone describe a contradiction between the nominal safety envelope and what operations in fact permitted. Recommendations, directives, training revisions, and design changes attempt to change the admissibility of future operations. A stabilized-approach policy, an improved checklist, a revised deicing procedure, or an airworthiness directive is not a comment on the past only. It is a transformation of what counts as a licit way for the system to continue.

\section{Why denominator problems do not kill the ledger}
None of this eliminates the value of denominators. Aviation certainly uses rates where it can: accidents per departure, per flight hour, by fleet type, or by operation class. Those measures are indispensable for some regulatory and managerial questions. But the field also knows that many actionable safety questions are misdescribed if forced into a single rate frame too early.

Suppose a set of reports concerns unstable approaches at a particular airport in unusual weather, involving a mixture of regional jets and turboprops, with similar last-minute workload spikes and phraseology confusions. One can want a denominator here --- approaches in comparable weather, by runway configuration, with the same traffic complexity --- and still rationally act before the estimate is fully stabilized. If the reports share severity markers and common contributing factors, if investigators or analysts can articulate the latent conditions, and if proposed interventions are modest relative to the downside risk, the absence of a polished global rate does not make inaction epistemically superior.

Indeed, a demand for only denominator-clean evidence can be perverse in safety practice. It privileges the questions easiest to count over the hazards easiest to understand. Near-miss systems exist partly because the catastrophic endpoint is too sparse, and too costly, to be the only admissible teacher.

This is a domain instance of a more general point about coherence and countability. Flyxion's \emph{Contradiction and Continuation} argues, in the Mutual Non-Determination Proposition of Part VI, that a system's logical or coding status and its structural coherence do not determine one another: a ledger can be well-formed and action-guiding while its statistical coding remains unsettled, just as a coding scheme can be tidy while the underlying safety picture stays incoherent. Denominator uncertainty is therefore not, by itself, evidence against the ledger's coherence.\footnote{Flyxion, \emph{Contradiction and Continuation}, Part VI, \url{https://github.com/standardgalactic/alphabet/tree/core/soup/book-latex-source}.}

Aviation's own quantitative risk literature makes the same point in a more formal idiom. Arnold Barnett describes a safety-analysis world increasingly concerned with precursor events and probabilistic updating rather than with waiting for crash counts alone.\footnote{Arnold Barnett, ``Aviation safety: a whole new world?'' \emph{Transportation Science} 44, no. 3 (2010): 289--302.} That is a more statistically rigorous version of the present essay's core claim: observable near-miss and precursor rates can rationally change one's estimate of catastrophic risk even when there is no single clean, global, i.i.d. denominator for the whole domain. In effect, mature aviation safety has formal ways to connect imperfect incident streams to rare-outcome risk without demanding that every admissible inference first become a raw population rate.

\section{Disconfirming instance: the limits are real}
A genuine counterweight comes from aviation's own self-critique. ASRS repeatedly warns that its data are ``soft data,'' affected by self-reporting bias and often uncorroborated beyond the reporter's perception. The public documentation states plainly that the existence of records on a topic cannot by itself establish prevalence in the National Airspace System. Underreporting remains a structural worry even in a trusted, non-punitive system. Some events are never seen, some are seen but not reported, and some are reported by the kinds of actors already most disposed toward conscientious disclosure.

There is also a substantive literature criticizing overly simple uses of the Swiss cheese model. Scholars such as Sidney Dekker and Erik Hollnagel have argued that the model can become too linear, too static, and too barrier-centric if treated as a complete theory of accidents rather than as a heuristic.\footnote{See, e.g., Sidney Dekker, \textit{Drift into Failure} (2011); Erik Hollnagel, \textit{Barriers and Accident Prevention} (2004).} Real socio-technical systems adapt, drift, and reorganize in ways that a stack-of-defenses picture can underdescribe. Those criticisms matter here because the present essay relies on the model only for a modest point: aviation explicitly treats latent, recurrent, organizational causes as safety-relevant. It does not rely on Swiss cheese as the final word on causation.

So aviation is not a proof that every incident ledger works. It is a proof that a ledger can remain useful despite imperfect denominators, voluntary near-miss data, and non-independence --- provided the field is disciplined about what claims the ledger may and may not support.

\section{Non-triviality conditions}
The central move of this essay would be wrong or empty under at least three conditions.

First, it would fail if aviation in fact made no practical use of such ledgers absent rate-quality denominator data. But it plainly does: ASRS alerting, \textit{CALLBACK}, coded trend review, and recommendation tracking all operate in advance of idealized exposure models.

Second, it would fail if the field's existing account already reduced everything to ordinary rate estimation and independent-event statistics. It does not. Aviation safety has long worked with human-factors coding, latent-condition analysis, and recommendation-based learning that are not reducible to raw counts.

Third, it would fail if a deliberate search for counterexamples showed that the apparent successes were merely rhetorical. Yet the strongest counterexamples found here --- underreporting, soft-data limitations, and criticism of simplified accident metaphors --- qualify the practice without defeating it. They narrow the lesson rather than erase it.

\section{A tentative why}
This section is explicitly speculative. Perhaps incident-ledger skepticism recurs because many fields inherit an image of knowledge from idealized science rather than from safety engineering. On that image, evidence is respectable when it looks like clean measurement over well-defined populations. Aviation developed under harsher constraints. Rare catastrophes are too sparse to teach fast enough; front-line operators know things formal audits miss; and the system must keep running while being corrected. Under those conditions, the relevant question shifts from ``is this an unbiased census?'' to ``is this an honest enough signal to alter the next loop of operation?''

That is close to the narrow sense in which \textit{Contradiction and Continuation}\footnote{Flyxion, \emph{Contradiction and Continuation} (draft), a study of Graham Priest's non-classical logics reworked as regimes of one constrained-representation theory, source at \url{https://github.com/standardgalactic/alphabet/blob/core/soup/book-latex-source/main.tex}.} treats consequence as preservation under an admissibility regime: one need not know everything globally before deciding whether a proposed next step can be licensed. The resemblance should not be overstated. But it helps explain why aviation's ledger is not primarily a book of counts. It is a disciplined filter for deciding which disturbances deserve to reshape future action.

\section{Conclusion}
Aviation safety reporting does not vindicate every proposed incident ledger. It does something more precise and more useful: it refutes the claim that denominator uncertainty, voluntary near-miss reporting, and non-independence make the ledger form itself epistemically bankrupt. ASRS and NTSB practice show that a mature field can extract durable safety knowledge from imperfect reports by coding them, constraining their use, merging perspectives, searching for latent common causes, and tying findings to \repair. The lesson for AI and other emerging domains is not that statistics do not matter. It is that the first mistake is often to ask a safety ledger to be a perfect rate table before allowing it to become a working instrument of learning.

\end{document}
